\subsubsection{SWORD}

Метод SWORD \cite{sword} в отличие от ранее рассмотренного метода Schism, который использует детализированную модель совместного доступа между вершинами графа, SWORD ориентирован на снижение вычислительной сложности за счёт использования сжатого представления рабочей нагрузки и инкрементального обновления разбиения.

Ключевое отличие SWORD заключается в том, что он не стремится к построению полной модели взаимодействий между всеми вершинами графа. Вместо этого используется приближённая структура, отражающая наиболее значимые зависимости, возникающие в реальной рабочей нагрузке. Это позволяет существенно повысить масштабируемость метода при сохранении приемлемого качества разбиения.

\textit{Общая концепция метода}

SWORD основан на идее сокращения пространства анализа рабочей нагрузки. Вместо построения плотного графа совместного доступа, как в Schism, используется разреженное представление взаимодействий между вершинами.

Основная гипотеза метода заключается в том, что для эффективного разбиения графа достаточно учитывать только наиболее значимые связи, возникающие в рабочей нагрузке. Менее частые или случайные взаимодействия могут быть отброшены без существенного ухудшения качества разбиения.

Таким образом, SWORD можно интерпретировать как метод, который заменяет точное моделирование рабочей нагрузки на её структурно сжатую аппроксимацию.

\textit{Модель рабочей нагрузки}

На первом этапе SWORD строит приближённую модель рабочей нагрузки на основе анализа журнала запросов $Q$. В отличие от Schism, где рассматриваются все пары совместного использования вершин, здесь применяется агрегирование информации.

Вводится функция значимости взаимодействия между вершинами:
\begin{equation*}
w(u, v) \approx f(Q),
\end{equation*}
где функция $f(Q)$ определяется на основе частоты совместного участия вершин $u$ и $v$ в запросах, однако вычисляется не через полный перебор всех пар, а через агрегированные структуры.

Далее формируется разреженный граф рабочей нагрузки:
\begin{equation*}
G_w = (V, E_w),
\end{equation*}
где множество рёбер $E_w$ определяется условием:
\begin{equation*}
E_w = \{(u, v) \mid w(u, v) > \theta\}.
\end{equation*}

Порог $\theta$ используется для отсечения слабых связей, которые не оказывают существенного влияния на стоимость выполнения запросов.

\textit{Сжатие и агрегация запросов}

Одним из ключевых этапов SWORD является предварительное сжатие рабочей нагрузки. Вместо анализа каждого запроса отдельно выполняется группировка запросов по схожести их структуры.

Пусть множество запросов $Q$ разбивается на кластеры:
\begin{equation*}
Q = \bigcup_{i=1}^{l} Q_i,
\end{equation*}
где каждый кластер $Q_i$ соответствует определённому типу доступа к графу.

В отличие от Schism, где точная статистика запросов используется напрямую, SWORD опирается на агрегированные характеристики кластеров. Для каждого кластера вычисляется обобщённая модель взаимодействий, которая используется для обновления графа $G_w$.

Такой подход позволяет существенно снизить вычислительные затраты при построении модели рабочей нагрузки, особенно при большом объёме запросов.

\textit{Инициализация разбиения}

После построения приближённой модели рабочей нагрузки выполняется инициализация разбиения графа. Начальное разбиение $\Pi^{(0)}$ формируется эвристически:
\begin{equation*}
\Pi^{(0)} = \{P_1^{(0)}, P_2^{(0)}, \dots, P_k^{(0)}\}.
\end{equation*}

В отличие от Schism, где качество начальной инициализации может существенно влиять на итоговый результат из-за более сложной модели оптимизации, SWORD менее чувствителен к начальному разбиению благодаря использованию локальных итеративных обновлений.

\textit{Функция стоимости разбиения}

Стоимость выполнения рабочей нагрузки определяется количеством межпартиционных взаимодействий, возникающих при обработке запросов.

Для запроса $q$ стоимость может быть выражена как:
\begin{equation*}
C(q) = g(\{P_i \mid v \in q, v \in P_i\}),
\end{equation*}
где функция $g$ отражает число переходов между партициями при выполнении запроса.

Общая стоимость системы определяется как:
\begin{equation*}
C(Q) = \sum_{q \in Q} C(q).
\end{equation*}

В отличие от Schism, где данная функция вычисляется на основе точного графа совместного доступа, в SWORD она является аппроксимацией, зависящей от разреженной модели $G_w$.

\textit{Итеративная локальная оптимизация}

Основной этап алгоритма SWORD заключается в итеративной оптимизации разбиения. Для каждой вершины $v \in V$ рассматривается возможность её перемещения между партициями.

Перемещение вершины из $P_i$ в $P_j$ принимается, если выполняется условие:
\begin{equation*}
\Delta C < 0.
\end{equation*}
Здесь $\Delta C$ обозначает изменение стоимости рабочей нагрузки после переноса вершины.

Процесс оптимизации выполняется локально, без глобального пересчёта всей модели взаимодействий. Это является важным отличием от Schism, где оптимизация тесно связана с точной структурой графа совместного доступа.

Локальный характер обновлений позволяет существенно снизить вычислительную сложность и повысить масштабируемость метода.

\textit{Инкрементальное обновление модели}

SWORD поддерживает механизм инкрементального обновления рабочей нагрузки. При поступлении новых запросов не требуется полное перестроение графа $G_w$. Вместо этого обновляются только те части модели, которые затронуты изменениями в рабочей нагрузке.

Это достигается за счёт хранения агрегированных статистик, позволяющих локально корректировать веса рёбер без глобального пересчёта всех взаимодействий.

Данный механизм особенно важен в условиях динамически изменяющейся нагрузки, характерной для реальных графовых приложений.

\textit{Балансировка партиций}

Как и в большинстве методов разбиения, в SWORD учитывается ограничение на размер партиций:
\begin{equation*}
|P_i| \leq (1 + \epsilon)\frac{|V|}{k}.
\end{equation*}
При перемещении вершин между партициями контролируется выполнение данного ограничения. В случае его нарушения выполняются корректирующие операции перераспределения.

\textit{Вычислительная сложность}

Основное преимущество SWORD заключается в снижении вычислительной сложности по сравнению с Schism. Построение модели рабочей нагрузки осуществляется на основе агрегированных данных, что уменьшает затраты на анализ запросов.

Итеративная оптимизация имеет следующую оценку сложности:
\begin{equation*}
O(I \cdot |V| \cdot k),
\end{equation*}
где $I$ обозначает число итераций до сходимости.

При этом ключевая экономия достигается на этапе построения графа рабочей нагрузки, который в SWORD не требует полного перебора всех пар вершин.

\textit{Сравнение со Schism}

Сопоставляя SWORD с Schism, можно выделить фундаментальное различие в уровне детализации модели рабочей нагрузки. Schism использует точную модель взаимодействий между вершинами графа, основанную на полном анализе запросов, тогда как SWORD применяет приближённую и разреженную модель.

Это приводит к следующему компромиссу: снижение точности моделирования компенсируется значительным увеличением масштабируемости. В результате SWORD оказывается более пригодным для систем с большим объёмом данных и высокой динамикой рабочей нагрузки.

Кроме того, SWORD лучше адаптируется к изменениям характера запросов благодаря инкрементальному обновлению модели, тогда как Schism требует более глобального пересчёта разбиения при существенных изменениях workload.

\textit{Характеристика получаемого разбиения}

Получаемое разбиение характеризуется локализацией наиболее частых взаимодействий между вершинами графа. Вероятность межпартиционных обращений снижается за счёт группировки часто совместно используемых данных.

При этом менее значимые взаимодействия могут оставаться разрезанными между партициями, что является следствием использования разреженной модели рабочей нагрузки.

Таким образом, итоговое разбиение представляет собой компромисс между точностью учёта взаимодействий и вычислительной эффективностью.